\chapter*{Abstract}
Standard Model (SM) electroweak theory (EW) and radiative corrections to the Weinberg angle ($\sin^2\theta_W$) predict that neutrinos, while electrically neutral, possess quantifiable electromagnetic properties. Specifically, one such property is the neutrino charge radius (NCR). We present how different experimental configurations involving reactor antineutrinos interacting via Coherent Elastic Neutrino-Nucleus Scattering (CE$\nu$NS) in liquid argon (LAr) can yield differing results depending on the values of $\sin^2\theta_W$ and the NCR, thus quantifying the sensitivity of such experiments to these SM parameters. The results show the projected sensitivity of reactor-based experiments near Mexico City to determining the Weinberg angle and the electron-neutrino NCR. We find that the projected precision on both parameters scales close to $1/\sqrt{N}$ with exposure time and improves by nearly an order of magnitude across the detector masses considered, from 10 to 200 kg, while the inclusion of a realistic quenching factor for argon, rather than an idealized unit response, degrades the sensitivity to both parameters by a factor of two to three due to the loss of low-energy recoil events below threshold. Even under this realistic scenario, a 200 kg detector exposed for 200 days is projected to constrain $\sin^2\theta_W$ to the $1$-$2\%$ level and the electron-neutrino charge radius to $\langle r_{\nu_e}^2\rangle \sim 2\times10^{-33}\ \text{cm}^2$, a precision competitive with existing low-energy determinations of the weak mixing angle and a concrete, testable target for the electron-neutrino electromagnetic structure using a source and detector technology already available at a national laboratory scale.
